\section{Conclusions and Policy Implications}
\label{sec:conclusions}

\subsection{Key Scientific Conclusions}
\label{subsec:concl_scientific}

This study integrated 43,937 NASA FIRMS satellite observations, official BBLSDLP (2019) peatland classifications, metric canal vectors in UTM Zone 49S, ECMWF ERA5-Land reanalysis, and cadastral concession layers to analyze wildfire drivers, predictive modeling, and governance in Central Kalimantan. The findings show that:
\begin{enumerate}
    \item Official SiPongi KLHK records (2018--2026) indicate that over 85\% of cumulative burned area occurred during two El Ni\~no episodes (2019 and 2023), confirming that macro-climatic drought controls fire extent.
    \item Spatial overlays demonstrated that 75.42\% of active fires occurred on unallocated Non-Concession land (\textit{Area Penggunaan Lain}, APL), whereas oil palm (HGU) and forestry (PBPH) concessions accounted for 16.56\% and 8.02\%, respectively. Fire mitigation is constrained by the absence of formal administration across unmanaged state peatlands.
    \item Peatlands accounted for 74.3\% of fire detections, with over 52.5\% located on peat domes deeper than 2\,m. In the ex-PLG landscape, primary canal distance was the second strongest predictor of severe Fire Radiative Power ($\text{mean } |\text{SHAP}| = 0.0322$).
    \item Coordinate-dependent Random Forest models (Model~A) collapsed under Spatial Block Cross-Validation (mean ROC-AUC fell from 0.7956 to 0.5399, and $F_1$-score was 0.0000 in four of five folds). In contrast, the coordinate-free biophysical model (Model~B) maintained positive severe fire detection across all five folds (spatial CV mean $F_1 = 0.1399$; spatial ROC-AUC = 0.4950), demonstrating that removing coordinates prevents spatial memorization collapse.
\end{enumerate}

\subsection{Actionable Policy Recommendations}
\label{subsec:concl_policy}

Based on these findings, we outline four policy interventions for disaster management authorities:
\begin{enumerate}
    \item Provincial disaster agencies should establish an emergency protocol elevating preparedness whenever meteorological stations record 12--14 consecutive rainless days ($<1.0$\,mm daily precipitation) in southern peat regencies, mobilizing patrols and suppression equipment before underground peat combustion begins.
    \item The Peatland and Mangrove Restoration Agency (BRGM) and the Ministry of Public Works should prioritize canal-blocking investments (box dams and compacted peat dams) within a 3,000-meter corridor along primary drainage canals in ex-PLG Blocks A and B to maintain phreatic water levels above $-40$\,cm.
    \item The Ministry of Environment and Forestry and regional governments should accelerate the assignment of abandoned ex-PLG APL peatlands into Social Forestry schemes, Forest Management Units (KPH), or customary stewardship, providing formal tenure and fire-prevention resources.
    \item Environmental law enforcement agencies should deploy automated spatial dashboards cross-referencing FIRMS thermal anomalies against corporate HGU and PBPH boundaries to monitor compliance and enforce liability under Law No.~32/2009.
\end{enumerate}

\subsection{Future Research}
\label{subsec:concl_future}

Future research can extend this work in three directions: (1) using Sentinel-1 InSAR interferometry to monitor surface peat subsidence as an indicator of groundwater table depletion; (2) deploying in-situ Internet-of-Things (IoT) groundwater level sensor networks across high-risk peat domes \citep{Dadap2021NatComm}; and (3) coupling machine learning fire intensity outputs with atmospheric chemical transport models (such as WRF-Chem) to forecast downwind $\text{PM}_{2.5}$ aerosol dispersion and transboundary smoke-haze exposure across Southeast Asia \citep{Kiely2021ACP}.
